\clearpage
\begin{proofbox}
\textbf{证明一｜最容易懂：从三格数表看出规律}
\begin{center}\begin{tikzpicture}[font=\small,>=stealth,line join=round]
\draw[gray!40,rounded corners=3pt] (0,0) rectangle (14.2,8.3);\node at (7.1,7.700000000000001) {例1：每走2，高度在3与负3之间交替};
\begin{scope}[shift={(1.5444444444444445,3.828947368421052)},x=0.5555555555555555cm,y=0.4842105263157894cm]\draw[->,gray] (-0.8,0)--(10,0) node[right]{$x$};\draw[->,gray] (0,-4.5)--(0,5) node[above]{$y$};\draw[blue,thick,solid,domain=-0.4:9.5,samples=180] plot(\x,{3*sin(deg(pi*\x/2))});\fill[orange!85!black] (1,3) circle(2.2pt);\fill[orange!85!black] (3,-3) circle(2.2pt);\fill[orange!85!black] (5,3) circle(2.2pt);\fill[orange!85!black] (7,-3) circle(2.2pt);\fill[orange!85!black] (9,3) circle(2.2pt);\foreach \v in {1,3,5,7,9}{\node[below=5pt] at (\v,0){$\v$};}\end{scope}\node[anchor=west,align=left,text width=5.0cm] at (8.4,4.0) {$f=3\sin\frac{\pi x}2$\\[9pt]$1,3,5,7,9$\\[9pt]对应 $3,-3,3,-3,3$\\[9pt]隔4回原值\\[9pt]$f(9)=f(1)=3$};
\node[text width=12.6cm,align=center] at (7.1,.65) {示例满足初值3；不能把每次相反误认为每次相等。};
\end{tikzpicture}\end{center}



\textbf{为什么这样想？} 只盯住某一个输入 $x$，别急着计算整个函数。设这一格的函数值是 $v$：
\[
\begin{array}{c|ccc}
\text{输入}&x&x+a&x+2a\\\hline
\text{输出}&v&-v&v
\end{array}
\]
第二格是第一格的相反数，这是题目直接告诉我们的；第三格又是第二格的相反数，于是 $-(-v)=v$。

\textbf{为什么能推广到所有 $x$？} 因为题目说的是“\textbf{每一个 $x$ 都这样}”，而不是只说某个特定数字。于是任取 $x$，输入相差 $2a$ 的两格函数值都相等，故 $f(x+2a)=f(x)$。

\medskip
\end{proofbox}\clearpage\begin{proofbox}\textbf{证明二｜考试最快：只换一次字母（推荐答题用）}
\begin{center}\begin{tikzpicture}[font=\small,>=stealth,line join=round]
\draw[gray!40,rounded corners=3pt] (0,0) rectangle (14.2,8.3);\node at (7.1,7.700000000000001) {隔2变为相反高度，隔4才回到原高度};
\begin{scope}[shift={(1.7153846153846155,3.8842857142857143)},x=0.7692307692307693cm,y=1.3142857142857143cm]\draw[->,gray] (-0.8,0)--(7,0) node[right]{$x$};\draw[->,gray] (0,-1.7)--(0,1.8) node[above]{$y$};\draw[blue,thick,solid,domain=-0.4:6.6,samples=180] plot(\x,{sin(deg(pi*\x/2))});\fill[orange!85!black] (1,1) circle(2.2pt);\fill[orange!85!black] (3,-1) circle(2.2pt);\fill[orange!85!black] (5,1) circle(2.2pt);\foreach \v/\h in {1/1,3/-1,5/1}{\draw[gray,dashed] (\v,0)--(\v,\h);\node[below=5pt] at (\v,0){$\v$};}\end{scope}\node[anchor=west,align=left,text width=5.0cm] at (8.4,4.0) {$f=\sin\frac{\pi x}2,\ a=2$\\[9pt]$f(1)=1$\\[9pt]$f(3)=-1$\\[9pt]$f(5)=1$\\[9pt]正周期可取4};
\node[text width=12.6cm,align=center] at (7.1,.65) {蓝线是满足条件的一条示例曲线；证明只用任意输入的两次变号。};
\end{tikzpicture}\end{center}



已知对任意 $x$，有 $f(x+a)=-f(x)$。

\textbf{第 1 步：}把条件中的 $x$ 换成 $x+a$。这是合法的，因为 $x+a$ 也是一个实数：
\[
f((x+a)+a)=-f(x+a).
\]
\textbf{第 2 步：}整理左边，再把原来的 $f(x+a)=-f(x)$ 代入右边：
\[
\boxed{\ f(x+2a)=-f(x+a)=-[-f(x)]=f(x).\ }
\]
\textbf{第 3 步：}若 $a>0$，正周期取 $2a$；若 $a<0$，因 $f(x+2a)=f(x)$ 也能推出 $f(x-2a)=f(x)$，故正周期取 $-2a$。统一得到 $2|a|$ 是一个正周期。

\medskip
\end{proofbox}\clearpage\begin{proofbox}\textbf{证明三｜从图像上的三个点看：横向走两步}
\begin{center}\begin{tikzpicture}[font=\small,>=stealth,line join=round]
\draw[gray!40,rounded corners=3pt] (0,0) rectangle (14.2,8.3);\node at (7.1,7.700000000000001) {证明三：从任意起点出发，同一过程处处成立};
\begin{scope}[shift={(1.7153846153846155,3.8842857142857143)},x=0.7692307692307693cm,y=1.3142857142857143cm]\draw[->,gray] (-0.8,0)--(7,0) node[right]{$x$};\draw[->,gray] (0,-1.7)--(0,1.8) node[above]{$y$};\draw[blue,thick,solid,domain=-0.4:6.6,samples=180] plot(\x,{sin(deg(pi*\x/2))});\fill[orange!85!black] (0.5,0.7071067811865476) circle(2.2pt);\fill[orange!85!black] (2.5,-0.7071067811865476) circle(2.2pt);\fill[orange!85!black] (4.5,0.7071067811865476) circle(2.2pt);\draw[orange!85!black,dashed] (.5,.70710678)--(4.5,.70710678);\node[above=5pt] at (.5,.70710678){$P$};\node[below=5pt] at (2.5,-.70710678){$Q$};\node[above=5pt] at (4.5,.70710678){$R$};\end{scope}\node[anchor=west,align=left,text width=5.0cm] at (8.4,4.0) {图取 $x=\frac12$\\[9pt]$P=(x,y)$\\[9pt]$Q=(x+2,-y)$\\[9pt]$R=(x+4,y)$\\[9pt]$P,R$ 横向相距4};
\node[text width=12.6cm,align=center] at (7.1,.65) {示例中任意改变起点，前后两次反号的关系仍然成立。};
\end{tikzpicture}\end{center}



从图像上\textbf{任意找一点} $P=(x,y)$，也就是说 $y=f(x)$。
\begin{enumerate}[leftmargin=*,itemsep=0.25em]
\item 输入加 $a$，函数值变号，所以 $Q=(x+a,-y)$ 也在图像上。
\item 从 $Q$ 的横坐标再加 $a$，函数值再变号，得到 $R=(x+2a,y)$，仍在图像上。
\item 点 $P$ 与点 $R$ 的高度相同，横坐标相隔 $2a$。\textbf{每一个图像上的点都能这样找到对应的点}，因此整幅图像每隔 $2|a|$ 会重复出现。
\end{enumerate}
此法只用点的坐标，特别适合向初中生解释“为什么是两次，而不是一次”。

\medskip
\end{proofbox}\clearpage\begin{proofbox}
\begin{center}\begin{tikzpicture}[font=\small,>=stealth,line join=round]
\draw[gray!40,rounded corners=3pt] (0,0) rectangle (14.2,8.3);\node at (7.1,7.700000000000001) {正反方向都一样：偶数步同高，奇数步相反};
\begin{scope}[shift={(3.266666666666667,3.8842857142857143)},x=0.8333333333333334cm,y=1.3142857142857143cm]\draw[->,gray] (-2.6,0)--(4.6,0) node[right]{$x$};\draw[->,gray] (0,-1.7)--(0,1.8) node[above]{$y$};\draw[blue,thick,solid,domain=-2.3:4.3,samples=180] plot(\x,{cos(deg(pi*\x))});\fill[orange!85!black] (-2,1.0) circle(2.2pt);\fill[orange!85!black] (-1,-1.0) circle(2.2pt);\fill[orange!85!black] (0,1) circle(2.2pt);\fill[orange!85!black] (1,-1) circle(2.2pt);\fill[orange!85!black] (2,1) circle(2.2pt);\fill[orange!85!black] (3,-1) circle(2.2pt);\fill[orange!85!black] (4,1) circle(2.2pt);\foreach \v in {-2,-1,1,2,3,4}{\ifnum\v=0\node[below left=5pt] at (\v,0){$\v$};\else\node[below=5pt] at (\v,0){$\v$};\fi}\end{scope}\node[anchor=west,align=left,text width=5.0cm] at (8.4,4.0) {图取 $a=1,\ f=\cos\pi x$\\[9pt]偶数输入：$f=1$\\[9pt]奇数输入：$f=-1$\\[9pt]负数步也满足\\[9pt]$f(x+na)=(-1)^nf(x)$};
\node[text width=12.6cm,align=center] at (7.1,.65) {各标数点来自同一周期曲线；整数步长可以向左或向右。};
\end{tikzpicture}\end{center}

\textbf{再多走一步｜为什么奇数次变号、偶数次不变？}

把刚才的过程不断接下去：$v,-v,v,-v,v,\dots$。每次只改变正负号，所以\textbf{经过偶数步}得到 $v$，\textbf{经过奇数步}得到 $-v$。往反方向走也一样，因为由 $f(x+a)=-f(x)$ 可推出 $f(x-a)=-f(x)$。所以对于任意整数 $n$，奇偶规律都成立。

\textbf{请区分：}这三种思路都证明了 $2|a|$ \textbf{是一个正周期}；任何一种都没有证明它\textbf{最小}，不能把“一个周期”擅自写成“最小正周期”。
\end{proofbox}
\clearpage
\begin{proofbox}
\textbf{补充证明｜把错位反号换成已经会的反号型}
\begin{center}\begin{tikzpicture}[font=\small,>=stealth,line join=round]
\draw[gray!40,rounded corners=3pt] (0,0) rectangle (14.2,8.3);
\node[text width=12.6cm,align=center] at (7.1,7.65) {错位反号：先认两个输入，再认两次移动};

\begin{scope}[shift={(1.8,4.3)},x=1.45cm,y=1.05cm]
\draw[->,gray] (-.4,0)--(6.8,0) node[right]{$u$};
\draw[->,gray] (0,-1.5)--(0,1.5) node[above left=3pt]{$y$};
\draw[blue,thick,domain=0:6.3,samples=160] plot(\x,{cos(deg(pi*\x/3))});
\foreach \u/\v/\lab in {0/1/t,3/-1/{t+d},6/1/{t+2d}}{
\draw[gray,dashed] (\u,0)--(\u,\v);\fill[orange!85!black] (\u,\v) circle(2.3pt);
\node[below=5pt] at (\u,-1.55) {$\lab$};}
\draw[<->,violet] (0,1.65)--(3,1.65) node[midway,above]{$d$};
\draw[<->,violet] (3,1.65)--(6,1.65) node[midway,above]{$d$};
\end{scope}
\node[align=center] at (7.1,1.3) {$t=x+b,\quad t+d=x+a,\quad d=a-b$};

\node[text width=12.6cm,align=center] at (7.1,.65) {图取 $d=3,t=0$，蓝线为 $y=\cos(\pi u/3)$；一般证明不依赖这条示例曲线。};
\end{tikzpicture}\end{center}


先取定义域为 $\mathbb R$。对任意 $t$，令 $t=x+b$，即在题设中取 $x=t-b$，得到
\[f(t+a-b)=-f(t).\]
设 $d=a-b\ne0$，就变成 $f(t+d)=-f(t)$。再把这个等式中的 $t$ 换成 $t+d$：
\[f(t+2d)=-f(t+d)=-[-f(t)]=f(t).\]
若 $d>0$，得到向右移动 $2d$ 后重复；若 $d<0$，将最后一式中的 $t$ 换成 $t-2d$，得 $f(t)=f(t-2d)$。所以统一取正数 $2|d|=2|a-b|$。
\textbf{这不是新的周期规律：}只是把原来的输入 $x+b$ 改叫 $t$，随后使用同样的两次变号。其他定义域按结论中的条件核对每次代入和来回平移。

\end{proofbox}
